\answer{ex:20:1}{$t_{\text{w}}=\tau_r\ln\frac{1-L}{1-H}=16.8$~h, $t_{\text{s}}=\tau_d\ln\frac HL=5.8$~h, period $22.5$~h; the daily swing of the thresholds pulls the oscillator to $24$~h: a phase-locked loop~\eqref{eq:pll}.}
\answer{ex:20:2}{(a)~$L=0.111$; with $L=0.092$ sleep would last $9.6$~h. (b)~By $22\%$ ($1/0.82=1.22$), not by $18\%$.}
\answer{ex:20:3}{$H=\frac{p-1}{p-1/q}=0.623$, $L=H/q=0.093$, where $p=e^{16/\tau_r}$, $q=e^{8/\tau_d}$. After being woken at $4$~h the phase shifts permanently (falling asleep $7.2$~h earlier); with swinging thresholds it returns within two to three days.}
\answer{ex:20:4}{(a)~$r_\pm=\frac12\bigl(1\pm\sqrt{1-4k/w}\bigr)=0.947,\ 0.053$; $a_{\text{hi}}=3.51$, $a_{\text{lo}}=0.49$. (b)~Activity $\approx0.33$~s, silence $\approx0.59$~s, period $\approx0.93$~s, active fraction $0.36$.}
\answer{ex:20:5}{$a_{\text{hi}}/r_+<b<a_{\text{lo}}/r_-$: $3.71<b<9.20$. \nt{ACET} lowers $b$ below $3.71$: the intersection moves to the upper branch, and the active state ($r\approx1$) is stable.}
\answer{ex:20:6}{$4.9$, $6.8$, $8.8$, $5.5$, $8.9$~h for $D=24$, $32$, $40$, $48$, $64$: the duration is set by the daily phase at which sleep begins, not by the debt.}
\answer{ex:20:7}{After B the error on A averages $0.51$ (from $0.19$ to $1.08$ over $20$ sets; a zero output gives $1$). Interleaved, both errors are below $10^{-4}$.}
\answer{ex:20:8}{Open problem. Uniform scaling preserves the weight ratios, but it preserves the responses of a threshold neuron only if the threshold is scaled by the same factor; interleaved replays change the ratios. That the large contacts stay unchanged contradicts purely uniform scaling.}
